%%==========================================================================
%% This is the very first version of a Streetwise LaTeX template. By
%% chosing the document class, it is easily convertible to a two-column IEEE
%% paper format, provided any conflicting packages and (new)commmands are
%% removed.
%%
%% Erwin de Gelder, January 6, 2018
%%==========================================================================

%%==========================================================================
%% Title:
%% Author(s):
%% To be published in:
%% Date of submission:
%% Date of submission R1:
%% Date of submission R2:
%% Date of final submission:
%%==========================================================================


%% Generic
\documentclass[10pt,final,a4paper,oneside,onecolumn]{article}

%% IEEE
%% Document class options: replace "draftcls" by "final" for final document.
%% All other options may just as well be omitted because they are the default values.
%\documentclass[10pt,final,journal,letterpaper,twoside,twocolumn]{IEEEtran}


%%==========================================================================
%% Document automation
%%==========================================================================

\def\reptitle{Completeness of real-world driving data}
\def\repauthor{Erwin de Gelder}


%%==========================================================================
%% Packages
%%==========================================================================

\usepackage[a4paper,left=3.5cm,right=3.5cm,top=3cm,bottom=3cm]{geometry} %% change page layout; remove for IEEE paper format
\usepackage[T1]{fontenc}                        %% output font encoding for international characters (e.g., accented)
\usepackage[cmex10]{amsmath}                    %% math typesetting; consider using the [cmex10] option
\usepackage{amssymb}                            %% special (symbol) fonts for math typesetting
\usepackage{amsthm}                             %% theorem styles
\usepackage{dsfont}                             %% double stroke roman fonts: the real numbers R: $\mathds{R}$
\usepackage{mathrsfs}                           %% formal script fonts: the Laplace transform L: $\mathscr{L}$
\usepackage[pdftex]{graphicx}                   %% graphics control; use dvips for TeXify; use pdftex for PDFTeXify
\usepackage{array}                              %% array functionality (array, tabular)
\usepackage{upgreek}                            %% upright Greek letters; add the prefix 'up', e.g. \upphi

\usepackage[utf8]{inputenc}   				 	%% utf8 support (required for biblatex)
\usepackage[style=ieee,doi=false,isbn=false,url=false,date=year,minbibnames=15,maxbibnames=15,backend=biber]{biblatex}
%\renewcommand*{\bibfont}{\footnotesize}		%% Use this for papers
\setlength{\biblabelsep}{\labelsep}
\bibliography{../bib}

\usepackage{stfloats}                           %% improved handling of floats
\usepackage{multirow}                           %% cells spanning multiple rows in tables
%\usepackage{subfigure}                         %% subfigures and corresponding captions (for use with IEEEconf.cls)
%\usepackage{subfig}                             %% subfigures (IEEEtran.cls: set caption=false)
\usepackage{fancyhdr}                           %% page headers and footers
\usepackage[official,left]{eurosym}             %% the euro symbol; command: \euro
\usepackage{appendix}                           %% appendix layout
\usepackage{xspace}                             %% add space after macro depending on context
\usepackage{verbatim}                           %% provides the comment environment
\usepackage[dutch,USenglish]{babel}             %% language support
\usepackage{wrapfig}                            %% wrapping text around figures
\usepackage{longtable}                          %% tables spanning multiple pages
\usepackage{pgfplots}                           %% support for TikZ figures (Matlab)
\usepackage[breaklinks=true,hidelinks,          %% implement hyperlinks (dvips yields minor problems with breaklinks;
            bookmarksnumbered=true]{hyperref}   %% IEEEtran: set bookmarks=false)
%\usepackage[hyphenbreaks]{breakurl}            %% allow line breaks in URLs (don't use with PDFTeX)
\usepackage{lmodern} 
\usepackage{etoolbox}							%% Needed for apptocmd later
\usepackage[capitalize]{cleveref}
\usepackage{units}
\usepackage{subcaption}
\usepackage{csquotes}							%% Quoted texts are typeset according to rules of main language

%%==========================================================================
%% Fancy headers and footers
%%==========================================================================

\newtoggle{standalone}
\toggletrue{standalone}
\pagestyle{fancy}                                       %% set page style
\fancyhf{}                                              %% clear all header & footer fields
\iftoggle{standalone}{%
	\fancyhead[L]{\includegraphics[width=15mm]{Streetwise}} %% define headers (LE: left field/even pages, etc.)
	\fancyhead[R]{\small\emph{\reptitle}}                   %% similar
	\fancyfoot[C]{\thepage}                                 %% define footer
	\setlength{\headheight}{18pt}                           %% increase head height to accommodate an oversized picture in the header
	\renewcommand*{\headrulewidth}{0.25pt}                  %% header line width
	%% Redefine the default "plain" page style (automatically activated by \maketitle, \section, ...)
	\fancypagestyle{plain}{%
		\fancyhf{}
		\fancyhead[C]{\includegraphics[width=30mm]{Streetwise}}
		\renewcommand{\headrulewidth}{0pt}
		\renewcommand{\footrulewidth}{0pt}
		\setlength{\headheight}{36pt}
	}
}{%
	\renewcommand*{\headrulewidth}{0pt}                  	%% No line in this case
	%% Redefine the default "plain" page style (automatically activated by \maketitle, \section, ...)
	\fancypagestyle{plain}{%
		\fancyhf{}
	}
}
\renewcommand*{\footrulewidth}{0pt}                     %% footer line width


%%==========================================================================
%% TikZ figures
%%==========================================================================

\newlength\figurewidth
\setlength\figurewidth{0.35\textwidth}              %% set figure width
\newlength\figureheight
\setlength\figureheight{0.3\textwidth}              %% set figure height
\pgfplotsset{every axis/.append style={
    scaled y ticks=false,
    scaled x ticks=false,
    y tick label style={/pgf/number format/fixed},
    x tick label style={/pgf/number format/fixed},
    legend style={font=\small}},
    compat=1.9}                                     %% PGFPlots package options
\usetikzlibrary{shapes.geometric, arrows, arrows.meta}
\usepgfplotslibrary{groupplots}
%\usetikzlibrary{external}                           %% Create pdf figures from TikZ. Use PDFTeXify ...
%\tikzexternalize[prefix=./tikz/]                    %% ... with --tex-option=--shell-escape switch.
%\tikzset{external/force remake}                    %% force pdf figure update


%%==========================================================================
%% User-defined commands
%%==========================================================================

\newcommand*{\mat}[1]{\mathbf{#1}}                              %% matrix/vector notation
\newcommand*{\matsym}[1]{\boldsymbol{#1}}                       %% matrix/vector notation for Greek letters
\newcommand*{\T}{^{\scriptscriptstyle\mathsf{T}}}               %% transpose operator
\newcommand*{\Hr}{^{\scriptscriptstyle\mathsf{H}}}              %% conjugate transpose operator
\newcommand*{\ud}{\mathrm{\,d}}                                 %% differential operator (upright d)
\newcommand*{\defeq}{\mathrel{\mathop:}=}                       %% definition sign :=
\newcommand*{\eqdef}{=\mathrel{\mathop:}}                       %% definition sign =:
\newcommand*{\ip}[2]{\left\langle#1\,{,}\,#2\right\rangle}      %% inner product
\newcommand*{\real}[1]{\mathrm{Re}(#1)}                         %% real part
\newcommand*{\imag}[1]{\mathrm{Im}(#1)}                         %% imaginary part
\newcommand*{\lsup}[1]{{}^{#1}\!}                               %% left superscript
\newcommand*{\hi}[1]{$^\text{#1}$}                              %% superscript in normal text
\newcommand*{\lo}[1]{$_\text{#1}$}                              %% subscript in normal text
\newcommand*{\w}[1]{\mathrm{#1}}                                %% multiple character super-/subscript in math mode
\newcommand*{\capskip}{\vspace{-12pt}}                          %% caption skip for figures with subfloats
\newcommand*{\etal}{et al.}                                     %% may be required for Natbib bibliography styles
\renewcommand*{\qedsymbol}{$\blacksquare$}                      %% redefine the end-of-proof symbol
\renewcommand*{\labelitemi}{$\bullet$}                          %% first level item list bullet
\renewcommand*{\labelitemii}{$-$}                               %% second level item list bullet
%\renewcommand*{\theenumi}{\textit{\roman{enumi}}}               %% first level enumerator
\renewcommand*{\labelenumi}{\theenumi.}
\renewcommand*{\theenumii}{\textit{\alph{enumii}}}              %% second level enumerator
\renewcommand*{\labelenumii}{\theenumii.}
\DeclareMathOperator{\tr}{tr}                                   %% trace of a matrix
\DeclareMathOperator{\sgn}{sgn}                                 %% signum function
\DeclareMathOperator{\atan}{atan}                               %% arc tangent

\newcommand{\deuclid}{d_{\textup{Euclidian}}}
\newcommand{\ddtw}{d_{\textup{DTW}}}
\newcommand{\imagi}{\textbf{j}}
\newcommand{\profile}[1]{\textbf{v}_{#1}}
\newcommand{\barprofile}[1]{\bar{\textbf{v}}_{#1}}
\newcommand{\parvstart}{v_{\textup{start}}}
\newcommand{\parvend}{v_{\textup{end}}}
\newcommand{\pardv}{v_{\Delta}}
\newcommand{\pardt}{t_{\textup{duration}}}
\newcommand{\expectation}[1]{\textup{E} \left[ #1 \right]}
\newcommand{\mise}[1]{\textup{MISE} \left( #1 \right)}
\newcommand{\entropy}[1]{\textup{H} \left[ #1 \right]}
\newcommand{\hatentropy}[1]{\hat{\textup{H}} \left[ #1 \right]}
\newcommand{\kl}[2]{\textup{KL} \left( #1 || #2 \right)}
\newcommand{\intinf}{\int_{-\infty}^{\infty}}

\crefname{figure}{Figure}{Figures}
\crefname{equation}{}{}
\Crefname{equation}{Equation}{Equations}

%%==========================================================================
%% User-defined environments
%%==========================================================================

\theoremstyle{plain}\newtheorem{definition}{Definition}[section]    %% definition
                    \newtheorem{theorem}{Theorem}[section]          %% theorem
                    \newtheorem{lemma}[theorem]{Lemma}              %% lemma
                    \newtheorem{corollary}[theorem]{Corollary}      %% corollary
                    \newtheorem{assumption}{Assumption}[section]    %% assumption
                    \newtheorem{condition}{Condition}[section]      %% condition
\theoremstyle{definition}\newtheorem{example}{Example}[section]     %% examples
\theoremstyle{remark}\newtheorem{remarkenv}{Remark}[section]        %% remarks
\newenvironment{remark}{\begin{remarkenv}}%
                       {\hfill$\blacklozenge$\end{remarkenv}}       %% end remark with a lozenge
\newenvironment{reviewer}{\itshape}{\upshape}                       %% environment for reviewer's comments


%%==========================================================================
%% Miscellaneous
%%==========================================================================

\graphicspath{{./../}{./figures/}{./more_figures/}} %% (graphicx) directory path for figures
%\setlength{\parindent}{0pt}                        %% no paragraph indentation
%\setlength{\parskip}{2ex}                          %% create empty line between paragraphs
%\interdisplaylinepenalty=2500                      %% (amsmath) allow for page breaks within multiline equations
%\numberwithin{equation}{section}                   %% (amsmath) include section number in equation numbering
%\numberwithin{figure}{section}                     %% (amsmath) include section number in figure numbering
%\numberwithin{table}{section}                      %% (amsmath) include section number in table numbering
\addtolength{\arraycolsep}{-0.5mm}                  %% squeeze matrix columns a little
\fnbelowfloat                                       %% (stfloats) put footnote below a float at the page bottom
\urlstyle{same}                                     %% (hyperref) use current font for URLs
\hypersetup{pdftitle={\reptitle},
            pdfauthor={\repauthor}}                 %% (hyperref) pdf properties title and author
%\raggedbottom                                      %% don't add inter-paragraph spacing to achieve \textheight
%\setlength\subfigcapskip{-3pt}                     %% (subfigure) distance between subfloat and subcaption
%\setlength\subfigbottomskip{-3pt}                  %% (subfigure) distance between subcaption and caption
%\renewcommand*{\subcapsize}{\small}                %% (subfigure) subcaption font size
\renewcommand*{\thesubfigure}{(\alph{subfigure})}   %% (subfig) implement 1(a) instead of 1a ...as subfigure reference
\captionsetup[subfloat]{labelformat=simple}        %% ... as subfigure reference
%\captionsetup[subfloat]{
%    farskip=0pt,
%    nearskip=8pt,
%    captionskip=1.5pt,
%    labelfont={small,bf},
%    textfont=small}                                %% (subfig) subfloat caption format
%\captionsetup[figure]{
%    labelfont={small,bf},
%    textfont=small}                                %% (subfig/caption) figure caption format
%\captionsetup[table]{
%    aboveskip=2pt,
%    labelfont={small,bf},
%    textfont={small,sc}}                           %% (subfig/caption) table caption format
\apptocmd{\thebibliography}{\raggedright}{}{}		%% Suppress badness warnings in bibliography


% Table stuff
\usepackage{booktabs}
\usepackage{tabularx}
\setlength{\heavyrulewidth}{0.1em}
\newcommand{\otoprule}{\midrule[\heavyrulewidth]}

%\usepackage{changebar}  %% When including the package earlier, it does not work for some reason...

%%==========================================================================
%% Begin document
%%==========================================================================

\begin{document}

\selectlanguage{USenglish}

\title{\textbf{\reptitle}}
\author{\repauthor}
\date{\today}
\maketitle

\tableofcontents

\newpage

\section{Introduction}
\label{sec:introduction}

% Scenario-based approach for assessment of AVs -> this scenarios are important
An important aspect in the development of automated vehicles (AVs) is the assessment of quality and performance aspects of the AVs, such as safety, comfort, and efficiency \cite{bengler2014threedecades, stellet2015taxonomy, wachenfeld2016release, putz2017pegasus, roesener2016scenariobased, kompass2015sicherheitsveranderung}. 
For legal and public acceptance, it is important that there is a clear definition of system performance and that there are quantitative measures for the system quality. 
The more traditional methods \cite{response2006code, ISO26262}, used for evaluation of driver assistance systems, are no longer valid for the assessment of quality and performance aspects of an AV \cite{wachenfeld2016release}. 
Therefore, a scenario-based approach is proposed \cite{roesener2016scenariobased, putz2017pegasus, kompass2015sicherheitsveranderung}. 
For the scenario-based assessment, proper specification of scenarios is crucial since they are directly reflected in test cases used for scenario-based assessment \cite{stellet2015taxonomy}.

% Completeness
To draw conclusions on how an automated vehicle would perform in real-world traffic, it is necessary to know how representative the scenario database, that is used for the scenario-based assessment of the automated vehicle, is. Therefore it is important to quantify how complete the scenario database is \cite{geyer2014, alvarez2017prospective, stellet2015taxonomy}.

% Structure of the document
In the next section, the problem that is addressed in this report is presented. In \cref{sec:method}, we show how this problem is approached. To illustrate this, an example is presented in \cref{sec:example}. Some questions are formulated in \cref{sec:questions}.

\section{Problem definition}
\label{sec:problem}

Before continuing our discussion about the scenario database, we will clarify the meaning of a scenario. A scenario is a quantitative description of the ego vehicle, its activities and/or goals, its dynamic environment (consisting of traffic environment and conditions) and its static environment. From the perspective of the ego vehicle, a scenario contains all relevant events. Scenarios can be categorized into so-called scenario classes. A scenario class refers to a qualitative description. \Cref{fig:scenario classes} shows the schematic representations of two examples of scenario classes.

It is assumed that the data can be described by a sequence of scenarios. Since these scenarios can categorized using different scenario classes, the problem of quantifying the completeness can now be divided into two subproblems.
\begin{enumerate}
	\item The first subproblem deals with the quantification of the completeness regarding the scenario classes.
	\item The second subproblem deals with the quantification of the completeness regarding the scenarios of a specific scenario class.
\end{enumerate}
In this report, we will only address the second subproblem. 

It is assumed that all scenarios that fall into the same scenario class can be parametrized in a similar manner using $d$ parameters. Let $n$ denote the number of scenarios for a specific scenario class. Consequently, we have $n$ parameter vectors, denoted by $x_i \in \mathbb{R}^d$ with $i\in\{1,...,n\}$, which describe the scenarios. To enable the generation of new scenarios from the same scenario class, we want to estimate the underlying distribution of $x_i$. Let $f(x)$ denote the true probability density function (pdf) and let $\hat{f}(x;n)$ denote the estimated pdf using $n$ parameter vectors. To quantify the completeness of the set of scenarios, represented by $x_i$, $i\in\{1,...,n\}$, we can make use of the Mean Integrated Squared Error (MISE), where the error is $f(x)-\hat{f}(x;n)$, i.e.,
\begin{equation} \label{eq:mise}
	\mise{n} = \expectation{ \int_{\mathbb{R}^d} \left(f(x) - \hat{f}(x;n)\right)^2 \ud x}.
\end{equation}

For the remainder of this report, it is assumed that the second subproblem can be addressed by estimating the MISE \cref{eq:mise} (see question \ref{question:mise}).

\begin{figure}
	\begin{subfigure}{\linewidth}
		\centering
		\includegraphics[scale=0.5]{"Scenario lead braking"}
		\caption{Schematic representation of the scenario class ``lead vehicle braking''.}
		\label{fig:lead vehicle braking}
	\end{subfigure}
	\begin{subfigure}{\linewidth}
		\centering
		\includegraphics[scale=0.5]{"Scenario cut-in"}
		\caption{Schematic representation of the scenario class ``cut-in''.}
		\label{fig:cut-in}
	\end{subfigure}
	\caption{Schematic representations of two different scenario classes. The blue vehicle represents the ego vehicle and the red vehicle represents a vehicle that belongs to the dynamic environment of the ego vehicle.}
	\label{fig:scenario classes}
\end{figure}



\section{Method}
\label{sec:method}

It is assumed that the pdf is estimated using Kernel Density Estimation (KDE) \cite{rosenblatt1956remarks, parzen1962estimation}, i.e.,
\begin{equation}
	\hat{f}(x;n) = \frac{1}{nh^d} \sum_{i=1}^n K \left( \frac{x - x_i}{h} \right),
\end{equation}
where $h$ denotes the bandwidth and $K(\cdot)$ denotes the kernel function. Furthermore, one-leave-out cross validation is used to compute the bandwidth $h$ because this minimizes the Kullback-Leibler divergence between the real pdf $f(x)$ and the estimated pdf $\hat{f}(x;n)$ \cite{turlach1993bandwidthselection}.

Chen \cite{chen2017tutorial} describes how the MISE \cref{eq:mise} can be estimated when KDE is used to estimate a pdf and when smoothness conditions apply:
\begin{equation} \label{eq:mise approximation}
	\mise{n} \approx \frac{h^4}{4} \sigma_K^4 \int_{\mathbb{R}^d} \left(\nabla^2 f(x)\right)^2 \, \textup{d}x + \frac{\mu_K}{nh^d}.
\end{equation}
Here, $\sigma_K^2 = \int_{\mathbb{R}^d} ||x||^2 K(x) \ud x$ and $\mu_K = \int_{\mathbb{R}^d} K(x)^2 \ud x$ are constants that depend on the choice of the kernel function $K(\cdot)$. In \cref{eq:mise approximation}, $\nabla^2 f(x)$ denotes the Laplacian of $f(x)$, i.e., $\nabla^2 f(x) = \sum_{i=1}^d \frac{\partial^2 f(x)}{\partial x_l^2}$ (this is similar to the trace of the Hessian of $f(x)$). The right-hand side of \cref{eq:mise approximation} is also know as the Asymptotic Mean Integrated Squared Error (AMISE) because it is similar to the MISE \cref{eq:mise} in the limit of $nh \rightarrow \infty$ and $h \rightarrow 0$.

The first term of the right-hand side of \cref{eq:mise approximation} corresponds to the bias introduced by smoothing the pdf. Therefore, this term approaches zero when $h \rightarrow 0$. However, when $h \rightarrow 0$, the variance goes to infinity, as can be seen by the second term of the right-hand side of \cref{eq:mise approximation}. Therefore, the optimal bandwidth is such that $h \rightarrow 0$ while $nh^d \rightarrow \infty$. Obviously, this is only achieved by having infinite data ($n \rightarrow \infty$).

We are unable to compute the approximation of the MISE, i.e., the right-hand side of \cref{eq:mise approximation}, because it involves the Laplacian of the true pdf $f(x)$. This term can be approximated, however, by using the estimated pdf $\hat{f}(x;n)$ \cite{calonico2018effect}. As a result, the MISE can be approximated as follows:
\begin{equation} \label{eq:mise final approximation}
	\mise{n} \approx \frac{h^4}{4} \sigma_K^4 \int_{\mathbb{R}^d} \left(\nabla^2 \hat{f}(x;n)\right)^2 \, \textup{d}x + \frac{\mu_K}{nh^d}.
\end{equation}



\section{Example}
\label{sec:example}

To illustrate the use of the approximated MISE \cref{eq:mise final approximation}, an example is presented. \Cref{fig:pdf} shows the pdf $f(x)$ that is used to generate data. The pdf $f(x)$ corresponds to a mixture of two normal distributions (with means of $-1$ and $1$ and standard deviations $0.5$ and $0.3$). 

\setlength\figurewidth{0.6\linewidth}
\setlength\figureheight{0.25\linewidth}
\begin{figure}
	\centering
	% This file was created by matplotlib2tikz v0.6.17.
\begin{tikzpicture}

\begin{axis}[
xlabel={$x$},
ylabel={Probability density},
xmin=-3, xmax=3,
ymin=-0.0184681719232084, ymax=0.388011272041166,
width=\figurewidth,
height=\figureheight,
tick align=outside,
tick pos=left,
x grid style={white!69.01960784313725!black},
y grid style={white!69.01960784313725!black},
axis x line*=bottom,
axis y line*=left,
every x tick/.style={black},
every y tick/.style={black}
]
\addplot [line width=0.7000000000000001pt, black, forget plot]
table {%
-3 0.0051667463394783
-2.99 0.00537706752410664
-2.98 0.00559483111163499
-2.97 0.00582024964842621
-2.96 0.0060535395815452
-2.95 0.00629492123931308
-2.94 0.00654461880712224
-2.93 0.00680286029835144
-2.92 0.00706987752022158
-2.91 0.00734590603443287
-2.9 0.00763118511242609
-2.89 0.00792595768511126
-2.88 0.00823047028690945
-2.87 0.00854497299395527
-2.86 0.00886971935631018
-2.85 0.00920496632403909
-2.84 0.00955097416700667
-2.83 0.00990800638825221
-2.82 0.0102763296308067
-2.81 0.0106562135778192
-2.8 0.0110479308458645
-2.79 0.0114517568713094
-2.78 0.0118679697896193
-2.77 0.0122968503074931
-2.76 0.0127386815677213
-2.75 0.0131937490066668
-2.74 0.0136623402042766
-2.73 0.0141447447265392
-2.72 0.0146412539603095
-2.71 0.0151521609404334
-2.7 0.0156777601691089
-2.69 0.0162183474274347
-2.68 0.0167742195791008
-2.67 0.01734567436619
-2.66 0.0179330101970674
-2.65 0.0185365259263441
-2.64 0.0191565206269158
-2.63 0.0197932933540838
-2.62 0.0204471429017821
-2.61 0.0211183675509413
-2.6 0.0218072648100373
-2.59 0.0225141311478814
-2.58 0.023239261718724
-2.57 0.0239829500797554
-2.56 0.0247454879011027
-2.55 0.025527164668434
-2.54 0.026328267378296
-2.53 0.0271490802263258
-2.52 0.0279898842884903
-2.51 0.0288509571955256
-2.5 0.0297325728007576
-2.49 0.030635000841507
-2.48 0.0315585065942914
-2.47 0.0325033505240556
-2.46 0.0334697879276774
-2.45 0.0344580685720086
-2.44 0.0354684363267296
-2.43 0.0365011287923104
-2.42 0.0375563769233859
-2.41 0.0386344046478707
-2.4 0.0397354284821519
-2.39 0.0408596571427155
-2.38 0.042007291154577
-2.37 0.0431785224569014
-2.36 0.0443735340062134
-2.35 0.0455924993776121
-2.34 0.0468355823644219
-2.33 0.0481029365767201
-2.32 0.0493947050392033
-2.31 0.0507110197888591
-2.3 0.0520520014729315
-2.29 0.0534177589476739
-2.28 0.0548083888784006
-2.27 0.0562239753413574
-2.26 0.0576645894279419
-2.25 0.0591302888518187
-2.24 0.0606211175594798
-2.23 0.0621371053448137
-2.22 0.0636782674682525
-2.21 0.0652446042810777
-2.2 0.0668361008554689
-2.19 0.0684527266208889
-2.18 0.0700944350074029
-2.17 0.0717611630965357
-2.16 0.0734528312802731
-2.15 0.0751693429288189
-2.14 0.0769105840677187
-2.13 0.078676423064967
-2.12 0.0804667103287082
-2.11 0.082281278016149
-2.1 0.0841199397542916
-2.09 0.0859824903730993
-2.08 0.0878687056516978
-2.07 0.0897783420782168
-2.06 0.0917111366238634
-2.05 0.0936668065318194
-2.04 0.0956450491215403
-2.03 0.09764554160903
-2.02 0.09966794094365
-2.01 0.101711883662016
-2 0.103776985759515
-1.99 0.10586284257998
-1.98 0.107969028724005
-1.97 0.110095097976434
-1.96 0.112240583253461
-1.95 0.114404996569825
-1.94 0.116587829026542
-1.93 0.11878855081957
-1.92 0.121006611269839
-1.91 0.123241438874997
-1.9 0.125492441383242
-1.89 0.127759005889563
-1.88 0.130040498954696
-1.87 0.132336266747074
-1.86 0.134645635208032
-1.85 0.136967910240477
-1.84 0.139302377921226
-1.83 0.141648304737188
-1.82 0.144004937845512
-1.81 0.146371505357815
-1.8 0.148747216648567
-1.79 0.15113126268767
-1.78 0.153522816397249
-1.77 0.15592103303262
-1.76 0.158325050587399
-1.75 0.160733990222639
-1.74 0.163146956719893
-1.73 0.165563038958021
-1.72 0.167981310413579
-1.71 0.170400829684527
-1.7 0.172820641037023
-1.69 0.175239774974984
-1.68 0.177657248832099
-1.67 0.180072067385914
-1.66 0.18248322349359
-1.65 0.184889698748903
-1.64 0.187290464160012
-1.63 0.189684480847488
-1.62 0.192070700762076
-1.61 0.194448067421612
-1.6 0.196815516666496
-1.59 0.199171977433089
-1.58 0.201516372544369
-1.57 0.203847619517157
-1.56 0.206164631385174
-1.55 0.208466317537208
-1.54 0.210751584569573
-1.53 0.2130193371521
-1.52 0.215268478906788
-1.51 0.217497913298296
-1.5 0.219706544535379
-1.49 0.221893278482381
-1.48 0.224057023579864
-1.47 0.226196691773434
-1.46 0.228311199449823
-1.45 0.230399468379241
-1.44 0.232460426663028
-1.43 0.234493009685602
-1.42 0.236496161069704
-1.41 0.238468833633914
-1.4 0.240409990351413
-1.39 0.242318605308979
-1.38 0.244193664665147
-1.37 0.246034167606532
-1.36 0.247839127301242
-1.35 0.249607571848372
-1.34 0.251338545222514
-1.33 0.253031108212268
-1.32 0.254684339351727
-1.31 0.256297335843899
-1.3 0.257869214475078
-1.29 0.259399112519145
-1.28 0.260886188630823
-1.27 0.262329623726905
-1.26 0.263728621854501
-1.25 0.265082411045364
-1.24 0.266390244155368
-1.23 0.267651399688237
-1.22 0.268865182602652
-1.21 0.270030925101868
-1.2 0.271147987405017
-1.19 0.272215758499275
-1.18 0.273233656872138
-1.17 0.274201131223024
-1.16 0.2751176611535
-1.15 0.275982757835427
-1.14 0.276795964656363
-1.13 0.277556857841597
-1.12 0.278265047052205
-1.11 0.278920175958573
-1.1 0.27952192278884
-1.09 0.280070000851776
-1.08 0.28056415903362
-1.07 0.281004182268453
-1.06 0.28138989198171
-1.05 0.28172114650646
-1.04 0.281997841472151
-1.03 0.282219910165499
-1.02 0.282387323863288
-1.01 0.282500092136846
-1 0.282558263128002
-0.99 0.282561923796383
-0.98 0.282511200137908
-0.97 0.282406257374403
-0.96 0.282247300114256
-0.95 0.282034572484095
-0.94 0.281768358231472
-0.93 0.281448980798587
-0.92 0.281076803367088
-0.91 0.280652228874041
-0.9 0.280175699999137
-0.89 0.279647699123292
-0.88 0.27906874825875
-0.87 0.278439408950863
-0.86 0.277760282151728
-0.85 0.277032008065857
-0.84 0.2762552659681
-0.83 0.275430773994042
-0.82 0.274559288903079
-0.81 0.27364160581444
-0.8 0.27267855791638
-0.79 0.271671016148788
-0.78 0.270619888859475
-0.77 0.269526121434386
-0.76 0.268390695901981
-0.75 0.26721463051204
-0.74 0.265998979289133
-0.73 0.264744831560988
-0.72 0.263453311461989
-0.71 0.262125577412021
-0.7 0.260762821570866
-0.69 0.259366269268347
-0.68 0.257937178410404
-0.67 0.256476838861257
-0.66 0.254986571801805
-0.65 0.253467729064402
-0.64 0.251921692444103
-0.63 0.250349872986483
-0.62 0.248753710252095
-0.61 0.247134671557616
-0.6 0.245494251193712
-0.59 0.243833969619621
-0.58 0.242155372634451
-0.57 0.24046003052514
-0.56 0.238749537191035
-0.55 0.237025509245004
-0.54 0.235289585090982
-0.53 0.233543423977833
-0.52 0.231788705029397
-0.51 0.230027126250564
-0.5 0.22826040350921
-0.49 0.226490269493812
-0.48 0.224718472646555
-0.47 0.222946776071716
-0.46 0.221176956419135
-0.45 0.219410802742544
-0.44 0.217650115332544
-0.43 0.215896704524014
-0.42 0.214152389477757
-0.41 0.212418996936151
-0.4 0.210698359952648
-0.39 0.208992316594926
-0.38 0.207302708621538
-0.37 0.205631380131934
-0.36 0.203980176189738
-0.35 0.202350941419217
-0.34 0.200745518574889
-0.33 0.199165747084277
-0.32 0.197613461563854
-0.31 0.196090490308269
-0.3 0.194598653753013
-0.29 0.19313976291071
-0.28 0.191715617781329
-0.27 0.190328005736631
-0.26 0.188978699879285
-0.25 0.187669457377121
-0.24 0.186402017773098
-0.23 0.185178101271653
-0.22 0.183999407002159
-0.21 0.182867611260352
-0.2 0.181784365728658
-0.19 0.180751295676472
-0.18 0.179769998141526
-0.17 0.178842040093626
-0.16 0.177968956582118
-0.15 0.17715224886857
-0.14 0.176393382546292
-0.13 0.17569378564841
-0.12 0.175054846746348
-0.11 0.174477913040697
-0.1 0.173964288446563
-0.0899999999999999 0.173515231675629
-0.0800000000000001 0.173131954317273
-0.0699999999999998 0.172815618921225
-0.0600000000000001 0.17256733708436
-0.0499999999999998 0.172388167544363
-0.04 0.172279114283089
-0.0299999999999998 0.172241124642621
-0.02 0.172275087457073
-0.00999999999999979 0.172381831203365
0 0.172562122174252
0.0100000000000002 0.172816662677044
0.02 0.173146089261506
0.0300000000000002 0.173550970980567
0.04 0.174031807687516
0.0500000000000003 0.174589028373473
0.0600000000000001 0.17522298954899
0.0700000000000003 0.175933973673688
0.0800000000000001 0.176722187637931
0.0899999999999999 0.177587761300538
0.1 0.178530746086628
0.11 0.179551113649681
0.12 0.180648754601943
0.13 0.181823477317294
0.14 0.183075006810732
0.15 0.18440298369855
0.16 0.185806963243341
0.17 0.187286414487842
0.18 0.188840719481647
0.19 0.190469172604724
0.2 0.192170979991594
0.21 0.193945259059946
0.22 0.195791038147376
0.23 0.197707256259786
0.24 0.199692762934888
0.25 0.201746318224077
0.26 0.203866592795827
0.27 0.20605216816353
0.28 0.208301537040587
0.29 0.210613103825297
0.3 0.212985185217933
0.31 0.215416010972121
0.32 0.217903724782466
0.33 0.220446385310044
0.34 0.223041967347206
0.35 0.225688363122788
0.36 0.228383383748615
0.37 0.231124760807842
0.38 0.233910148085405
0.39 0.236737123440547
0.4 0.239603190821053
0.41 0.242505782418525
0.42 0.24544226096369
0.43 0.248409922160404
0.44 0.251405997256667
0.45 0.254427655750652
0.46 0.257472008229368
0.47 0.260536109337273
0.48 0.263616960871767
0.49 0.266711515002191
0.5 0.269816677608583
0.51 0.272929311736128
0.52 0.27604624116087
0.53 0.279164254061976
0.54 0.282280106795447
0.55 0.285390527763928
0.56 0.288492221376897
0.57 0.291581872095257
0.58 0.294656148554039
0.59 0.29771170775666
0.6 0.300745199333918
0.61 0.303753269860636
0.62 0.306732567222653
0.63 0.309679745026623
0.64 0.31259146704489
0.65 0.315464411687506
0.66 0.318295276493306
0.67 0.321080782631788
0.68 0.323817679407431
0.69 0.326502748757958
0.7 0.329132809737977
0.71 0.331704722979358
0.72 0.334215395119643
0.73 0.336661783189808
0.74 0.33904089895263
0.75 0.341349813183
0.76 0.343585659881512
0.77 0.345745640412768
0.78 0.347827027559912
0.79 0.349827169487013
0.8 0.351743493601077
0.81 0.353573510305612
0.82 0.355314816637859
0.83 0.356965099782002
0.84 0.358522140450913
0.85 0.359983816129219
0.86 0.361348104170754
0.87 0.362613084743753
0.88 0.363776943617449
0.89 0.36483797478406
0.9 0.365794582910501
0.91 0.366645285614522
0.92 0.367388715560331
0.93 0.368023622369191
0.94 0.368548874340829
0.95 0.368963459981977
0.96 0.369266489338738
0.97 0.369457195129952
0.98 0.369534933679149
0.99 0.369499185643156
1 0.369349556535883
1.01 0.36908577704626
1.02 0.368707703149793
1.03 0.368215316013645
1.04 0.367608721695664
1.05 0.366888150638201
1.06 0.366053956958062
1.07 0.365106617534389
1.08 0.364046730896733
1.09 0.362875015916025
1.1 0.361592310301591
1.11 0.360199568907819
1.12 0.358697861854478
1.13 0.357088372465121
1.14 0.355372395028413
1.15 0.353551332387578
1.16 0.351626693363578
1.17 0.349600090017947
1.18 0.347473234761571
1.19 0.345247937316015
1.2 0.342926101534303
1.21 0.340509722088321
1.22 0.338000881030314
1.23 0.33540174423613
1.24 0.332714557738149
1.25 0.329941643955967
1.26 0.327085397833125
1.27 0.3241482828883
1.28 0.321132827189508
1.29 0.318041619259972
1.3 0.314877303924392
1.31 0.311642578104395
1.32 0.308340186572005
1.33 0.304972917669944
1.34 0.301543599007593
1.35 0.298055093141387
1.36 0.294510293248371
1.37 0.290912118801534
1.38 0.287263511255486
1.39 0.283567429750848
1.4 0.279826846845651
1.41 0.276044744281818
1.42 0.272224108794642
1.43 0.268367927972985
1.44 0.264479186177656
1.45 0.260560860525242
1.46 0.256615916944371
1.47 0.252647306311142
1.48 0.248657960670161
1.49 0.244650789547333
1.5 0.240628676360264
1.51 0.236594474931774
1.52 0.232551006111749
1.53 0.228501054512173
1.54 0.224447365359873
1.55 0.220392641471148
1.56 0.216339540352097
1.57 0.212290671428111
1.58 0.208248593405631
1.59 0.204215811768912
1.6 0.200194776414164
1.61 0.1961878794231
1.62 0.192197452977527
1.63 0.188225767416296
1.64 0.184275029435542
1.65 0.180347380432821
1.66 0.176444894995382
1.67 0.172569579532496
1.68 0.168723371051431
1.69 0.164908136076318
1.7 0.161125669708879
1.71 0.157377694829646
1.72 0.153665861438039
1.73 0.149991746129361
1.74 0.146356851706521
1.75 0.142762606924028
1.76 0.139210366361552
1.77 0.13570141042411
1.78 0.132236945465741
1.79 0.128818104033289
1.8 0.125445945226761
1.81 0.122121455172518
1.82 0.118845547605422
1.83 0.115619064555874
1.84 0.112442777137601
1.85 0.109317386431871
1.86 0.106243524463756
1.87 0.103221755265929
1.88 0.100252576025463
1.89 0.0973364183089753
1.9 0.0944736493614468
1.91 0.091664573474017
1.92 0.0889094334160046
1.93 0.0862084119264255
1.94 0.0835616332602638
1.95 0.0809691647847806
1.96 0.0784310186211632
1.97 0.0759471533268698
1.98 0.073517475614063
1.99 0.0711418420995978
2 0.0688200610820904
2.01 0.0665518943416869
2.02 0.0643370589582297
2.03 0.0621752291436253
2.04 0.0600660380843174
2.05 0.0580090797898876
2.06 0.0560039109439214
2.07 0.054050052753406
2.08 0.0521469927930545
2.09 0.0502941868410912
2.1 0.0484910607031715
2.11 0.0467370120212529
2.12 0.0450314120643826
2.13 0.0433736074985203
2.14 0.0417629221326598
2.15 0.0401986586386762
2.16 0.038680100242474
2.17 0.0372065123841689
2.18 0.0357771443451939
2.19 0.0343912308403751
2.2 0.0330479935731743
2.21 0.0317466427524541
2.22 0.0304863785692694
2.23 0.029266392632342
2.24 0.028085869361019
2.25 0.0269439873346649
2.26 0.0258399205975763
2.27 0.024772839918647
2.28 0.0237419140051493
2.29 0.0227463106701216
2.3 0.0217851979529895
2.31 0.0208577451931602
2.32 0.0199631240564569
2.33 0.019100509514366
2.34 0.0182690807761867
2.35 0.0174680221742701
2.36 0.0166965240026385
2.37 0.0159537833093662
2.38 0.0152390046431955
2.39 0.0145514007549439
2.4 0.0138901932543356
2.41 0.0132546132229675
2.42 0.0126439017841843
2.43 0.0120573106307051
2.44 0.0114941025108965
2.45 0.0109535516746468
2.46 0.0104349442798379
2.47 0.00993757876045978
2.48 0.00946076615744805
2.49 0.00900383041336154
2.5 0.00856610863204305
2.51 0.0081469513044355
2.52 0.00774572250174325
2.53 0.00736180003714709
2.54 0.00699457559729306
2.55 0.00664345484478508
2.56 0.00630785749291598
2.57 0.00598721735387412
2.58 0.00568098236166017
2.59 0.00538861457094565
2.6 0.0051095901330957
2.61 0.00484339925056927
2.62 0.00458954611089613
2.63 0.00434754880141539
2.64 0.00411693920594205
2.65 0.00389726288450818
2.66 0.00368807893730415
2.67 0.00348895985392085
2.68 0.00329949134896977
2.69 0.00311927218513029
2.7 0.00294791398464616
2.71 0.00278504103026373
2.72 0.00263029005657488
2.73 0.00248331003269625
2.74 0.00234376193718502
2.75 0.00221131852605901
2.76 0.0020856640947567
2.77 0.00196649423483898
2.78 0.00185351558620208
2.79 0.0017464455855368
2.8 0.00164501221173685
2.81 0.0015489537289249
2.82 0.00145801842773232
2.83 0.00137196436543572
2.84 0.00129055910552075
2.85 0.00121357945721161
2.86 0.00114081121547332
2.87 0.00107204890196223
2.88 0.00100709550737067
2.89 0.00094576223558088
2.9 0.00088786825001527
2.91 0.00083324042254097
2.92 0.000781713085259585
2.93 0.000733127785486239
2.94 0.000687333044196061
2.95 0.000644184118191438
2.96 0.000603542766219095
2.97 0.00056527701924307
2.98 0.000529260955057211
2.99 0.000495374477399596
3 0.000463503099710759
};
\addplot [line width=0.7000000000000001pt, gray, dashed, forget plot]
table {%
-3 8.1664388086521e-06
-2.99 8.86176776665391e-06
-2.98 9.6133975447717e-06
-2.97 1.04256233288926e-05
-2.96 1.13030451130026e-05
-2.95 1.22505874362978e-05
-2.94 1.32735202490426e-05
-2.93 1.43774809605026e-05
-2.92 1.55684977239872e-05
-2.91 1.68530140157314e-05
-2.9 1.82379145660389e-05
-2.89 1.97305527027708e-05
-2.88 2.13387791689166e-05
-2.87 2.30709724775906e-05
-2.86 2.49360708693774e-05
-2.85 2.69436059384712e-05
-2.84 2.91037379955335e-05
-2.83 3.1427293236591e-05
-2.82 3.39258027886395e-05
-2.81 3.6611543703857e-05
-2.8 3.94975819754935e-05
-2.79 4.25978176495442e-05
-2.78 4.59270321072488e-05
-2.77 4.95009375942538e-05
-2.76 5.33362290729395e-05
-2.75 5.74506384749045e-05
-2.74 6.18629914309427e-05
-2.73 6.65932665559876e-05
-2.72 7.16626573664484e-05
-2.71 7.7093636907102e-05
-2.7 8.29100251642051e-05
-2.69 8.91370593407499e-05
-2.68 9.58014670687875e-05
-2.67 0.000102931542632461
-2.66 0.000110557226273808
-2.65 0.000118710186651502
-2.64 0.000127423906520495
-2.63 0.000136733771697913
-2.62 0.000146677163377671
-2.61 0.000157293553852907
-2.6 0.000168624605701617
-2.59 0.00018071427448675
-2.58 0.000193608915017404
-2.57 0.000207357391212732
-2.56 0.000222011189604652
-2.55 0.000237624536509459
-2.54 0.00025425451889187
-2.53 0.000271961208938083
-2.52 0.000290807792346734
-2.51 0.000310860700338639
-2.5 0.000332189745377368
-2.49 0.000354868260583548
-2.48 0.000378973242815819
-2.47 0.000404585499380953
-2.46 0.000431789798324551
-2.45 0.000460675022242013
-2.44 0.000491334325537159
-2.43 0.000523865295042935
-2.42 0.000558370113905031
-2.41 0.000594955728615039
-2.4 0.000633734019064902
-2.39 0.000674821971478961
-2.38 0.000718341854063819
-2.37 0.000764421395199469
-2.36 0.000813193963977954
-2.35 0.00086479875287778
-2.34 0.00091938096234404
-2.33 0.000977091987024998
-2.32 0.00103808960339657
-2.31 0.00110253815848592
-2.3 0.0011706087593851
-2.29 0.00124247946322454
-2.28 0.00131833546725512
-2.27 0.00139836929866568
-2.26 0.00148278100374092
-2.25 0.00157177833594255
-2.24 0.00166557694247391
-2.23 0.00176440054886572
-2.22 0.00186848114109809
-2.21 0.00197805914475114
-2.2 0.00209338360065374
-2.19 0.00221471233647766
-2.18 0.00234231213370166
-2.17 0.00247645888934846
-2.16 0.00261743777187522
-2.15 0.00276554337057728
-2.14 0.00292107983784399
-2.13 0.00308436102358529
-2.12 0.0032557106011283
-2.11 0.00343546218386455
-2.1 0.00362395943191086
-2.09 0.0038215561480302
-2.08 0.00402861636204314
-2.07 0.00424551440294671
-2.06 0.00447263495794396
-2.05 0.00471037311757662
-2.04 0.00495913440614299
-2.03 0.00521933479657534
-2.02 0.00549140070894445
-2.01 0.0057757689917548
-2 0.0060728868851914
-1.99 0.00638321196547899
-1.98 0.00670721206951642
-1.97 0.00704536519895367
-1.96 0.00739815940288547
-1.95 0.00776609263834544
-1.94 0.00814967260779641
-1.93 0.00854941657282809
-1.92 0.00896585114329068
-1.91 0.00939951204111396
-1.9 0.0098509438380856
-1.89 0.0103206996668888
-1.88 0.0108093409047305
-1.87 0.0113174368289232
-1.86 0.0118455642438231
-1.85 0.0123943070785629
-1.84 0.0129642559550666
-1.83 0.013556007725875
-1.82 0.0141701649813649
-1.81 0.0148073355259967
-1.8 0.0154681318232826
-1.79 0.0161531704092285
-1.78 0.0168630712740655
-1.77 0.0175984572121554
-1.76 0.0183599531400235
-1.75 0.0191481853825463
-1.74 0.0199637809273986
-1.73 0.0208073666479424
-1.72 0.0216795684948254
-1.71 0.0225810106566381
-1.7 0.0235123146900695
-1.69 0.0244740986200898
-1.68 0.0254669760107816
-1.67 0.0264915550075337
-1.66 0.0275484373514117
-1.65 0.0286382173666116
-1.64 0.0297614809220083
-1.63 0.0309188043679047
-1.62 0.0321107534491946
-1.61 0.0333378821962483
-1.6 0.0346007317949365
-1.59 0.0358998294373079
-1.58 0.037235687154538
-1.57 0.0386088006338665
-1.56 0.0400196480213446
-1.55 0.0414686887123062
-1.54 0.0429563621315804
-1.53 0.0444830865055521
-1.52 0.0460492576282754
-1.51 0.0476552476239293
-1.5 0.0493014037079956
-1.49 0.0509880469496212
-1.48 0.0527154710377046
-1.47 0.0544839410533243
-1.46 0.0562936922511944
-1.45 0.0581449288529007
-1.44 0.0600378228547278
-1.43 0.061972512852946
-1.42 0.0639491028894692
-1.41 0.065967661320843
-1.4 0.0680282197135534
-1.39 0.0701307717686746
-1.38 0.0722752722788968
-1.37 0.0744616361209867
-1.36 0.0766897372867376
-1.35 0.0789594079554642
-1.34 0.0812704376110851
-1.33 0.0836225722068149
-1.32 0.0860155133804601
-1.31 0.0884489177232733
-1.3 0.0909223961052741
-1.29 0.0934355130598898
-1.28 0.0959877862306993
-1.27 0.0985786858829938
-1.26 0.101207634482779
-1.25 0.103874006345753
-1.24 0.106577127358689
-1.23 0.109316274775536
-1.22 0.112090677090447
-1.21 0.114899513989779
-1.2 0.117741916385013
-1.19 0.120616966528356
-1.18 0.123523698212654
-1.17 0.126461097057078
-1.16 0.129428100879865
-1.15 0.132423600159217
-1.14 0.135446438583287
-1.13 0.138495413689979
-1.12 0.141569277597086
-1.11 0.144666737823093
-1.1 0.147786458198762
-1.09 0.150927059869403
-1.08 0.1540871223875
-1.07 0.157265184895167
-1.06 0.160459747395669
-1.05 0.163669272113007
-1.04 0.166892184938355
-1.03 0.17012687696191
-1.02 0.173371706088462
-1.01 0.176624998734786
-1 0.179885051606734
-0.99 0.183150133553646
-0.98 0.186418487497521
-0.97 0.189688332434143
-0.96 0.192957865503134
-0.95 0.196225264123732
-0.94 0.199488688192849
-0.93 0.202746282341785
-0.92 0.205996178247795
-0.91 0.209236496996486
-0.9 0.212465351490892
-0.89 0.215680848902868
-0.88 0.218881093162329
-0.87 0.222064187479675
-0.86 0.225228236896635
-0.85 0.228371350860643
-0.84 0.231491645817723
-0.83 0.234587247818787
-0.82 0.23765629513416
-0.81 0.240696940871063
-0.8 0.243707355588743
-0.79 0.246685729905884
-0.78 0.249630277094908
-0.77 0.252539235657769
-0.76 0.255410871877826
-0.75 0.258243482342414
-0.74 0.26103539643075
-0.73 0.263784978761866
-0.72 0.266490631597324
-0.71 0.269150797193518
-0.7 0.271763960098506
-0.69 0.274328649388375
-0.68 0.276843440838287
-0.67 0.279306959023492
-0.66 0.281717879345709
-0.65 0.284074929980497
-0.64 0.286376893741348
-0.63 0.288622609856464
-0.62 0.290810975654381
-0.61 0.29294094815477
-0.6 0.295011545561037
-0.59 0.297021848651511
-0.58 0.298971002066283
-0.57 0.300858215487022
-0.56 0.302682764707317
-0.55 0.304443992591389
-0.54 0.306141309919284
-0.53 0.307774196116927
-0.52 0.309342199869707
-0.51 0.310844939618556
-0.5 0.312282103937755
-0.49 0.313653451794018
-0.48 0.314958812686675
-0.47 0.316198086669078
-0.46 0.317371244251649
-0.45 0.318478326187268
-0.44 0.319519443140001
-0.43 0.320494775238432
-0.42 0.321404571515163
-0.41 0.322249149234312
-0.4 0.323028893109089
-0.39 0.323744254411823
-0.38 0.324395749979032
-0.37 0.324983961114395
-0.36 0.325509532392705
-0.35 0.325973170368103
-0.34 0.326375642190121
-0.33 0.326717774131241
-0.32 0.327000450029882
-0.31 0.327224609652883
-0.3 0.327391246981737
-0.29 0.327501408426949
-0.28 0.327556190975045
-0.27 0.327556740272872
-0.26 0.327504248653923
-0.25 0.327399953111517
-0.24 0.327245133223743
-0.23 0.32704110903511
-0.22 0.326789238899917
-0.21 0.326490917292348
-0.2 0.326147572588347
-0.19 0.325760664824266
-0.18 0.325331683437325
-0.17 0.324862144992825
-0.16 0.324353590903029
-0.15 0.323807585142578
-0.14 0.323225711965173
-0.13 0.322609573626229
-0.12 0.32196078811603
-0.11 0.321280986907837
-0.1 0.320571812725259
-0.0899999999999999 0.319834917333023
-0.0800000000000001 0.319071959355165
-0.0699999999999998 0.318284602124456
-0.0600000000000001 0.317474511566724
-0.0499999999999998 0.316643354123535
-0.04 0.315792794716509
-0.0299999999999998 0.314924494756326
-0.02 0.314040110199294
-0.00999999999999979 0.313141289654113
0 0.312229672541257
0.0100000000000002 0.31130688730718
0.02 0.310374549695302
0.0300000000000002 0.309434261075528
0.04 0.308487606833803
0.0500000000000003 0.307536154822976
0.0600000000000001 0.306581453876023
0.0700000000000003 0.305625032382448
0.0800000000000001 0.304668396928432
0.0899999999999999 0.303713031001108
0.1 0.302760393757091
0.11 0.301811918855196
0.12 0.300869013353034
0.13 0.299933056667023
0.14 0.299005399595086
0.15 0.298087363401183
0.16 0.297180238960591
0.17 0.296285285964699
0.18 0.295403732183918
0.19 0.294536772787134
0.2 0.293685569716023
0.21 0.292851251112374
0.22 0.292034910796478
0.23 0.29123760779451
0.24 0.290460365912753
0.25 0.289704173356402
0.26 0.288969982390642
0.27 0.288258709041624
0.28 0.287571232834903
0.29 0.286908396568885
0.3 0.286271006120811
0.31 0.285659830282782
0.32 0.285075600625356
0.33 0.28451901138623
0.34 0.283990719381609
0.35 0.283491343937837
0.36 0.283021466840972
0.37 0.282581632302034
0.38 0.282172346935718
0.39 0.281794079750488
0.4 0.281447262148016
0.41 0.281132287930093
0.42 0.280849513311206
0.43 0.280599256935134
0.44 0.280381799894032
0.45 0.280197385748622
0.46 0.280046220548248
0.47 0.279928472849717
0.48 0.279844273733986
0.49 0.279793716819953
0.5 0.279776858274724
0.51 0.279793716819953
0.52 0.279844273733986
0.53 0.279928472849717
0.54 0.280046220548248
0.55 0.280197385748622
0.56 0.280381799894032
0.57 0.280599256935134
0.58 0.280849513311206
0.59 0.281132287930093
0.6 0.281447262148016
0.61 0.281794079750488
0.62 0.282172346935718
0.63 0.282581632302034
0.64 0.283021466840972
0.65 0.283491343937837
0.66 0.283990719381609
0.67 0.28451901138623
0.68 0.285075600625356
0.69 0.285659830282782
0.7 0.286271006120811
0.71 0.286908396568885
0.72 0.287571232834903
0.73 0.288258709041624
0.74 0.288969982390642
0.75 0.289704173356402
0.76 0.290460365912753
0.77 0.29123760779451
0.78 0.292034910796478
0.79 0.292851251112374
0.8 0.293685569716023
0.81 0.294536772787134
0.82 0.295403732183918
0.83 0.296285285964699
0.84 0.297180238960591
0.85 0.298087363401183
0.86 0.299005399595086
0.87 0.299933056667023
0.88 0.300869013353034
0.89 0.301811918855196
0.9 0.302760393757091
0.91 0.303713031001108
0.92 0.304668396928432
0.93 0.305625032382448
0.94 0.306581453876023
0.95 0.307536154822976
0.96 0.308487606833803
0.97 0.309434261075528
0.98 0.310374549695302
0.99 0.31130688730718
1 0.312229672541257
1.01 0.313141289654113
1.02 0.314040110199295
1.03 0.314924494756326
1.04 0.315792794716509
1.05 0.316643354123535
1.06 0.317474511566724
1.07 0.318284602124456
1.08 0.319071959355165
1.09 0.319834917333023
1.1 0.320571812725259
1.11 0.321280986907837
1.12 0.32196078811603
1.13 0.322609573626229
1.14 0.323225711965173
1.15 0.323807585142578
1.16 0.324353590903029
1.17 0.324862144992825
1.18 0.325331683437325
1.19 0.325760664824266
1.2 0.326147572588347
1.21 0.326490917292348
1.22 0.326789238899917
1.23 0.32704110903511
1.24 0.327245133223743
1.25 0.327399953111517
1.26 0.327504248653923
1.27 0.327556740272872
1.28 0.327556190975045
1.29 0.327501408426949
1.3 0.327391246981737
1.31 0.327224609652883
1.32 0.327000450029882
1.33 0.326717774131241
1.34 0.326375642190121
1.35 0.325973170368103
1.36 0.325509532392705
1.37 0.324983961114395
1.38 0.324395749979032
1.39 0.323744254411823
1.4 0.323028893109089
1.41 0.322249149234312
1.42 0.321404571515163
1.43 0.320494775238432
1.44 0.319519443140001
1.45 0.318478326187268
1.46 0.317371244251649
1.47 0.316198086669078
1.48 0.314958812686675
1.49 0.313653451794018
1.5 0.312282103937755
1.51 0.310844939618556
1.52 0.309342199869707
1.53 0.307774196116927
1.54 0.306141309919284
1.55 0.304443992591389
1.56 0.302682764707317
1.57 0.300858215487022
1.58 0.298971002066283
1.59 0.297021848651511
1.6 0.295011545561037
1.61 0.29294094815477
1.62 0.290810975654381
1.63 0.288622609856464
1.64 0.286376893741348
1.65 0.284074929980497
1.66 0.281717879345709
1.67 0.279306959023491
1.68 0.276843440838287
1.69 0.274328649388375
1.7 0.271763960098506
1.71 0.269150797193518
1.72 0.266490631597324
1.73 0.263784978761866
1.74 0.261035396430749
1.75 0.258243482342414
1.76 0.255410871877826
1.77 0.252539235657769
1.78 0.249630277094908
1.79 0.246685729905884
1.8 0.243707355588743
1.81 0.240696940871063
1.82 0.23765629513416
1.83 0.234587247818787
1.84 0.231491645817723
1.85 0.228371350860643
1.86 0.225228236896635
1.87 0.222064187479675
1.88 0.218881093162329
1.89 0.215680848902868
1.9 0.212465351490892
1.91 0.209236496996486
1.92 0.205996178247795
1.93 0.202746282341785
1.94 0.199488688192848
1.95 0.196225264123732
1.96 0.192957865503134
1.97 0.189688332434143
1.98 0.186418487497521
1.99 0.183150133553645
2 0.179885051606734
2.01 0.176624998734787
2.02 0.173371706088462
2.03 0.17012687696191
2.04 0.166892184938355
2.05 0.163669272113007
2.06 0.160459747395669
2.07 0.157265184895167
2.08 0.1540871223875
2.09 0.150927059869403
2.1 0.147786458198762
2.11 0.144666737823093
2.12 0.141569277597086
2.13 0.138495413689979
2.14 0.135446438583287
2.15 0.132423600159217
2.16 0.129428100879865
2.17 0.126461097057078
2.18 0.123523698212654
2.19 0.120616966528356
2.2 0.117741916385013
2.21 0.114899513989779
2.22 0.112090677090447
2.23 0.109316274775536
2.24 0.106577127358688
2.25 0.103874006345753
2.26 0.101207634482779
2.27 0.0985786858829937
2.28 0.0959877862306992
2.29 0.0934355130598898
2.3 0.0909223961052742
2.31 0.0884489177232732
2.32 0.0860155133804601
2.33 0.0836225722068149
2.34 0.0812704376110851
2.35 0.078959407955464
2.36 0.0766897372867375
2.37 0.0744616361209867
2.38 0.0722752722788968
2.39 0.0701307717686746
2.4 0.0680282197135533
2.41 0.065967661320843
2.42 0.0639491028894692
2.43 0.0619725128529461
2.44 0.0600378228547277
2.45 0.0581449288529006
2.46 0.0562936922511944
2.47 0.0544839410533244
2.48 0.0527154710377046
2.49 0.0509880469496212
2.5 0.0493014037079956
2.51 0.0476552476239293
2.52 0.0460492576282754
2.53 0.0444830865055521
2.54 0.0429563621315804
2.55 0.0414686887123063
2.56 0.0400196480213445
2.57 0.0386088006338665
2.58 0.037235687154538
2.59 0.035899829437308
2.6 0.0346007317949364
2.61 0.0333378821962483
2.62 0.0321107534491946
2.63 0.0309188043679047
2.64 0.0297614809220083
2.65 0.0286382173666116
2.66 0.0275484373514117
2.67 0.0264915550075337
2.68 0.0254669760107816
2.69 0.0244740986200898
2.7 0.0235123146900695
2.71 0.0225810106566381
2.72 0.0216795684948254
2.73 0.0208073666479423
2.74 0.0199637809273985
2.75 0.0191481853825463
2.76 0.0183599531400235
2.77 0.0175984572121554
2.78 0.0168630712740655
2.79 0.0161531704092285
2.8 0.0154681318232826
2.81 0.0148073355259967
2.82 0.0141701649813649
2.83 0.013556007725875
2.84 0.0129642559550666
2.85 0.0123943070785628
2.86 0.011845564243823
2.87 0.0113174368289232
2.88 0.0108093409047305
2.89 0.0103206996668889
2.9 0.00985094383808558
2.91 0.00939951204111395
2.92 0.00896585114329068
2.93 0.0085494165728281
2.94 0.00814967260779639
2.95 0.00776609263834543
2.96 0.00739815940288547
2.97 0.00704536519895367
2.98 0.0067072120695164
2.99 0.00638321196547898
3 0.0060728868851914
};
\end{axis}

\end{tikzpicture}
	\caption{The true distribution $f(x)$ that is used for the example. The pdf $f(x)$ corresponds to a mixture of two normal distributions (with means of $-1$ and $1$ and standard deviations $0.5$ and $0.3$). }
	\label{fig:pdf}
\end{figure}

To compare the approximated MISE of \cref{eq:mise final approximation} with the MISE of \cref{eq:mise}, the MISE needs to be computed. This is done by approximating the expectation in \cref{eq:mise} by taking the average of 100 experiments. The result is shown in \cref{fig:mise example}. The result shows that the approximated MISE of \cref{eq:mise final approximation} follows the same trend as the real MISE.

\setlength\figureheight{0.35\linewidth}
\begin{figure}
	\centering
	% This file was created by matplotlib2tikz v0.7.5.
\begin{tikzpicture}

\begin{axis}[
height=\figureheight,
log basis x={10},
log basis y={10},
tick align=outside,
tick pos=left,
width=\figurewidth,
x grid style={white!69.01960784313725!black},
xlabel={Number of samples},
xmajorgrids,
xmin=100, xmax=5000,
xmode=log,
xtick style={color=black},
y grid style={white!69.01960784313725!black},
ylabel={MISE},
ymajorgrids,
ymin=0.000398955915484946, ymax=0.0186585956806614,
ymode=log,
ytick style={color=black}
]
\path [draw=white!80.0!black, fill=white!80.0!black]
(axis cs:100,0.0156664988761491)
--(axis cs:100,0.00867859030734705)
--(axis cs:108,0.00770635037031353)
--(axis cs:117,0.00739957817229819)
--(axis cs:127,0.0060271782360017)
--(axis cs:138,0.00608378244738758)
--(axis cs:149,0.00710413508546981)
--(axis cs:161,0.0087789071781041)
--(axis cs:175,0.00782714536388698)
--(axis cs:189,0.00813143132236401)
--(axis cs:205,0.0061063293463563)
--(axis cs:222,0.00515627762915232)
--(axis cs:241,0.00455292320776853)
--(axis cs:261,0.00459052848555843)
--(axis cs:282,0.00418063030879111)
--(axis cs:306,0.00388765959817649)
--(axis cs:331,0.00383186441820095)
--(axis cs:359,0.00359673930613909)
--(axis cs:389,0.0033279770943303)
--(axis cs:421,0.00334800483523584)
--(axis cs:456,0.00295497036820648)
--(axis cs:494,0.0027941629102757)
--(axis cs:535,0.00277516246346377)
--(axis cs:579,0.00278994705570234)
--(axis cs:627,0.00252381470876108)
--(axis cs:679,0.00224900815839357)
--(axis cs:736,0.00211289247079706)
--(axis cs:797,0.00196884838018689)
--(axis cs:863,0.00189836691995214)
--(axis cs:935,0.00186617451835314)
--(axis cs:1013,0.00179180488112653)
--(axis cs:1097,0.00166348411868397)
--(axis cs:1188,0.00155119965503882)
--(axis cs:1287,0.00148869731555646)
--(axis cs:1394,0.00134091164379462)
--(axis cs:1510,0.00126125217181022)
--(axis cs:1635,0.00119515080555445)
--(axis cs:1771,0.00112660648405963)
--(axis cs:1918,0.00108931933865783)
--(axis cs:2078,0.000980838609676316)
--(axis cs:2250,0.000965286471555012)
--(axis cs:2437,0.000989817715414935)
--(axis cs:2640,0.000915629410011999)
--(axis cs:2859,0.000890564706293059)
--(axis cs:3097,0.000910705057971742)
--(axis cs:3354,0.000834590573932385)
--(axis cs:3633,0.000712279786946419)
--(axis cs:3935,0.000669028783329945)
--(axis cs:4262,0.000572772743529408)
--(axis cs:4616,0.000544199507793738)
--(axis cs:5000,0.000517560162858556)
--(axis cs:5000,0.000687026936708465)
--(axis cs:5000,0.000687026936708465)
--(axis cs:4616,0.000734639703176721)
--(axis cs:4262,0.000780258435914869)
--(axis cs:3935,0.000899063796778749)
--(axis cs:3633,0.000955801084398279)
--(axis cs:3354,0.00109786025766935)
--(axis cs:3097,0.00119538603258129)
--(axis cs:2859,0.00121118284463861)
--(axis cs:2640,0.00125772277615397)
--(axis cs:2437,0.00137334923298444)
--(axis cs:2250,0.00133450405044003)
--(axis cs:2078,0.00137272976278189)
--(axis cs:1918,0.00152590046034461)
--(axis cs:1771,0.00159568074021147)
--(axis cs:1635,0.00169004596013437)
--(axis cs:1510,0.00181482121145608)
--(axis cs:1394,0.00187581593942072)
--(axis cs:1287,0.00207210683079026)
--(axis cs:1188,0.00215719951611934)
--(axis cs:1097,0.00234525453873345)
--(axis cs:1013,0.00250095555874643)
--(axis cs:935,0.00263226648333643)
--(axis cs:863,0.00278077303218856)
--(axis cs:797,0.00296433004362662)
--(axis cs:736,0.0031762980537263)
--(axis cs:679,0.00346134792219142)
--(axis cs:627,0.00378394459695842)
--(axis cs:579,0.00412051964051918)
--(axis cs:535,0.0042393587279419)
--(axis cs:494,0.00443326910866449)
--(axis cs:456,0.00467169398639971)
--(axis cs:421,0.0050822046734161)
--(axis cs:389,0.00530164602220987)
--(axis cs:359,0.0056734398287967)
--(axis cs:331,0.0061118631333199)
--(axis cs:306,0.00644469702716999)
--(axis cs:282,0.00692515747729661)
--(axis cs:261,0.00756176090321755)
--(axis cs:241,0.00782636072710048)
--(axis cs:222,0.00857326361731711)
--(axis cs:205,0.00964317396175354)
--(axis cs:189,0.0122847892501114)
--(axis cs:175,0.0119683853862226)
--(axis cs:161,0.013203617263897)
--(axis cs:149,0.0118519287068655)
--(axis cs:138,0.0120226952113701)
--(axis cs:127,0.0128799905414387)
--(axis cs:117,0.0138437673850392)
--(axis cs:108,0.0144585748542336)
--(axis cs:100,0.0156664988761491)
--cycle;

\addplot [line width=2.0pt, black]
table {%
100 0.00859558582257646
108 0.00812354841097345
117 0.00767960999089539
127 0.00724156968333626
138 0.00672419333167824
149 0.00660572197053461
161 0.00769490075185661
175 0.0070536539874133
189 0.00734188669662077
205 0.00560004520248748
222 0.00497320029942687
241 0.00458341225138433
261 0.0044193967132435
282 0.00411277466976824
306 0.00391024863593416
331 0.00371502137121955
359 0.00345931539882796
389 0.00325192901866895
421 0.00303951306817579
456 0.00280851639425498
494 0.00271323547625052
535 0.00264215520746293
579 0.00258984470938745
627 0.00237308394789097
679 0.00219208629809438
736 0.00204851187007183
797 0.00192558480862347
863 0.00184092307760793
935 0.00174854822001538
1013 0.00167027944732962
1097 0.00154219545609572
1188 0.00141474033936746
1287 0.00135239133962665
1394 0.00123566339441748
1510 0.00116519558723847
1635 0.00109955350545929
1771 0.00105316800665945
1918 0.000998804164198836
2078 0.000938865893448118
2250 0.000909986663373976
2437 0.00092481765671545
2640 0.000852028795585114
2859 0.000824250734934949
3097 0.000833149925120803
3354 0.0007652585851753
3633 0.000661896203261313
3935 0.000612342178439652
4262 0.000532870354595377
4616 0.000499248115923246
5000 0.000475151288126954
};
\addplot [line width=2.0pt, gray, dashed]
table {%
100 0.0124222154744985
108 0.0114047924111372
117 0.0110856998560715
127 0.00908932470364904
138 0.00878628215427483
149 0.00986933776249793
161 0.0123917508129278
175 0.0105715440500389
189 0.0110696465088209
205 0.00858025418107767
222 0.00747320359552377
241 0.00632682919019152
261 0.00661365722756003
282 0.00596364492904691
306 0.00537243882058895
331 0.00518129038406335
359 0.00495020249335441
389 0.00456345424271336
421 0.00455613142317589
456 0.00414160123883459
494 0.00392694207708443
535 0.00389371221570144
579 0.00377971960642565
627 0.00337406939576007
679 0.00290753753157886
736 0.00266130816611476
797 0.00246875767172887
863 0.00238672019361078
935 0.00233647608913928
1013 0.00224908007641507
1097 0.00210606135141524
1188 0.00194369485769856
1287 0.00186906375094596
1394 0.00161076503782372
1510 0.00154713660477321
1635 0.00144645969912583
1771 0.00138043601579382
1918 0.00134242289195381
2078 0.00114594022552776
2250 0.00113333040382035
2437 0.00117663255862454
2640 0.00107084030741098
2859 0.00105339633754792
3097 0.00109031230531131
3354 0.000987151733632913
3633 0.000840833236739904
3935 0.000793096424147307
4262 0.000683616696267234
4616 0.00064411866574568
5000 0.000614216157862005
};
\end{axis}

\end{tikzpicture}
	\caption{Comparison of the MISE of \cref{eq:mise} (blue) and the approximation of the MISE using \cref{eq:mise final approximation} (orange) for different number of data points.}
	\label{fig:mise example}
\end{figure}

To illustrate that the result in \cref{fig:mise example} is not a mere coincidence as it depends on randomly generated data, the experiment is repeated 100 times. \Cref{fig:mise example repeats} shows the result of repeating the experiment 100 times. It shows that the average follows the same trend as the real MISE. Furthermore, the deviation is not very large.

\begin{figure}
	\centering
	% This file was created by matplotlib2tikz v0.6.17.
\begin{tikzpicture}

\definecolor{color0}{rgb}{0.12156862745098,0.466666666666667,0.705882352941177}
\definecolor{color1}{rgb}{1,0.498039215686275,0.0549019607843137}

\begin{axis}[
xlabel={Number of samples},
ylabel={MISE},
xmin=100, xmax=5000,
ymin=0.000400499833529209, ymax=0.0172049297359885,
xmode=log,
ymode=log,
width=\figurewidth,
height=\figureheight,
tick align=outside,
tick pos=left,
xmajorgrids,
x grid style={white!69.01960784313725!black},
ymajorgrids,
y grid style={white!69.01960784313725!black},
legend style={draw=white!80.0!black},
legend entries={{Real MISE},{Estimated MISE}},
legend cell align={left}
]
\addlegendimage{no markers, color0}
\addlegendimage{no markers, color1}
\addplot [line width=2.0pt, color0]
table {%
100 0.00859558582257646
108 0.00812354841097345
117 0.00767960999089539
127 0.00724156968333626
138 0.00672419333167824
149 0.00660572197053461
161 0.00769490075185661
175 0.0070536539874133
189 0.00734188669662077
205 0.00560004520248748
222 0.00497320029942687
241 0.00458341225138433
261 0.0044193967132435
282 0.00411277466976824
306 0.00391024863593416
331 0.00371502137121955
359 0.00345931539882796
389 0.00325192901866895
421 0.00303951306817579
456 0.00280851639425498
494 0.00271323547625052
535 0.00264215520746293
579 0.00258984470938745
627 0.00237308394789097
679 0.00219208629809438
736 0.00204851187007183
797 0.00192558480862347
863 0.00184092307760793
935 0.00174854822001538
1013 0.00167027944732962
1097 0.00154219545609572
1188 0.00141474033936746
1287 0.00135239133962665
1394 0.00123566339441748
1510 0.00116519558723847
1635 0.00109955350545929
1771 0.00105316800665945
1918 0.000998804164198836
2078 0.000938865893448118
2250 0.000909986663373976
2437 0.00092481765671545
2640 0.000852028795585114
2859 0.000824250734934949
3097 0.000833149925120803
3354 0.0007652585851753
3633 0.000661896203261313
3935 0.000612342178439652
4262 0.000532870354595377
4616 0.000499248115923246
5000 0.000475151288126954
};
\addplot [line width=2.0pt, color1]
table {%
100 0.0121725445917481
108 0.0110824626122735
117 0.0106216727786687
127 0.00945358438872021
138 0.00905323882937882
149 0.00947803189616764
161 0.0109912622210006
175 0.00989776537505481
189 0.0102081102862377
205 0.00787475165405492
222 0.00686477062323472
241 0.0061896419674345
261 0.00607614469438799
282 0.00555289389304386
306 0.00516617831267324
331 0.00497186377576043
359 0.00463508956746789
389 0.00431481155827009
421 0.00421510475432597
456 0.0038133321773031
494 0.0036137160094701
535 0.00350726059570284
579 0.00345523334811076
627 0.00315387965285975
679 0.0028551780402925
736 0.00264459526226168
797 0.00246658921190675
863 0.00233956997607035
935 0.00224922050084478
1013 0.00214638021993648
1097 0.00200436932870871
1188 0.00185419958557908
1287 0.00178040207317336
1394 0.00160836379160767
1510 0.00153803669163315
1635 0.00144259838284441
1771 0.00136114361213555
1918 0.00130760989950122
2078 0.0011767841862291
2250 0.00114989526099752
2437 0.00118158347419969
2640 0.00108667609308299
2859 0.00105087377546584
3097 0.00105304554527651
3354 0.00096622541580087
3633 0.00083404043567235
3935 0.000784046290054347
4262 0.000676515589722138
4616 0.00063941960548523
5000 0.00060229354978351
};
\addplot [line width=2.0pt, color1, dashed, forget plot]
table {%
100 0.0145018474480154
108 0.0133332041069136
117 0.0127697358495823
127 0.0117378551571992
138 0.011032876417373
149 0.0110606297699662
161 0.0124661655829315
175 0.0112781787158334
189 0.0115925629288202
205 0.00905369985918734
222 0.00800376595262298
241 0.00728078780721182
261 0.00706655550027436
282 0.00646773628254569
306 0.00601852412233774
331 0.00573186334746674
359 0.00532732307502043
389 0.00497270120089661
421 0.00479317136705272
456 0.00438557338336751
494 0.00416008474226636
535 0.00399532601719555
579 0.00389875754304971
627 0.00357392294892553
679 0.00325929129489178
736 0.00299906378990477
797 0.00279841643305333
863 0.00263370534681583
935 0.00250458448917255
1013 0.00238276377914312
1097 0.00223162613539187
1188 0.00205619953927258
1287 0.00197487191158463
1394 0.00178666522348304
1510 0.00172255970484844
1635 0.00160756343437105
1771 0.0015175016975195
1918 0.00145313694006348
2078 0.00130741457059762
2250 0.00127296778729253
2437 0.00130942731338952
2640 0.00120070721513031
2859 0.00115774648824769
3097 0.00114793920347969
3354 0.00105398197704653
3633 0.000915214201489636
3935 0.000860724627870615
4262 0.000745677487183958
4616 0.000702899670612891
5000 0.000658782474400146
};
\addplot [line width=2.0pt, color1, dashed, forget plot]
table {%
100 0.00984324173548073
108 0.00883172111763353
117 0.00847360970775502
127 0.0071693136202412
138 0.00707360124138466
149 0.00789543402236909
161 0.00951635885906959
175 0.00851735203427625
189 0.00882365764365524
205 0.00669580344892251
222 0.00572577529384646
241 0.00509849612765719
261 0.00508573388850162
282 0.00463805150354203
306 0.00431383250300874
331 0.00421186420405411
359 0.00394285605991536
389 0.00365692191564356
421 0.00363703814159922
456 0.00324109097123869
494 0.00306734727667383
535 0.00301919517421012
579 0.00301170915317182
627 0.00273383635679397
679 0.00245106478569321
736 0.0022901267346186
797 0.00213476199076017
863 0.00204543460532487
935 0.00199385651251702
1013 0.00190999666072985
1097 0.00177711252202555
1188 0.00165219963188557
1287 0.0015859322347621
1394 0.00143006235973231
1510 0.00135351367841786
1635 0.00127763333131777
1771 0.0012047855267516
1918 0.00116208285893896
2078 0.00104615380186058
2250 0.00102682273470252
2437 0.00105373963500985
2640 0.000972644971035661
2859 0.000944001062683984
3097 0.000958151887073333
3354 0.000878468854555213
3633 0.000752866669855063
3935 0.000707367952238079
4262 0.000607353692260317
4616 0.000575939540357569
5000 0.000545804625166874
};
\end{axis}

\end{tikzpicture}
	\caption{The real MISE is show by the blue line. The average of the approximation for 100 repeated experiments is shown by the solid orange line. The dashed orange lines represent the average of the approximation plus or minus twice the standard deviation.}
	\label{fig:mise example repeats}
\end{figure}

\section{Questions}
\label{sec:questions}

Based on the theory I described above, I have some questions:

\begin{enumerate}
	\item \label{question:mise} Does it seem correct to assume that the MISE can be used to quantify the completeness of the set of scenarios belonging to a specific scenario class? Do I overlook something?
	\item Given that the probability density function is estimated using Kernel Density Estimation, would there be other methods useful/applicable for estimating how reliable our estimated probability density function is?
	\item According to Chen \cite{chen2017tutorial}, the approximation of the MISE using the AMISE of \cref{eq:mise approximation} is valid if the smoothness conditions apply, i.e., the real pdf $f(x)$ needs to be smooth. That made me think that it is reasonable to approximate the MISE using \cref{eq:mise final approximation}. Is this too simplistic, i.e., do I overlook something?
	\item When searching on this topic, I found many papers that use the (A)MISE. However, if these papers use experiments at all, they use synthetic data (i.e., similar to what I did in the example). Do you know some studies where the same theory is applied for some real-life applications?
\end{enumerate}

\printbibliography

\end{document} 